\documentclass[a4paper,twocolumn,9pt]{bxjsarticle}

% ==========================================================
% 余白設定: 上下左右 約1インチ (25.4mm)
% ==========================================================
\usepackage[a4paper,margin=1in]{geometry}

% パッケージ
\usepackage{amsmath,amssymb}
\usepackage{url}
\usepackage{titlesec}
\usepackage{enumitem}

% セクション間隔の調整 (1ページに厳密に収めるための最適化)
\titlespacing*{\section}{0pt}{4pt plus 1pt minus 1pt}{2pt plus 1pt minus 1pt}
\titlespacing*{\subsection}{0pt}{3pt plus 1pt minus 1pt}{1pt plus 1pt minus 1pt}
\setlist{nosep,leftmargin=1.2em}

% 標題
\title{\vspace{-1.5em}\bfseries デュアルモード型ECプラットフォーム\\「WonderMarket」の設計と実装\vspace{-0.5em}}
\author{開発プロジェクトチーム}
\date{\vspace{-0.8em}2026年10月\vspace{-1.2em}}

\begin{document}

\maketitle

% ==========================================================
% 概要 (Abstract)
% ==========================================================
\begin{abstract}
\small
本稿では，Web標準技術（HTML5/CSS3/モダンJavaScript）およびPython軽量RESTful APIバックエンドを統合したハイブリッド型ECプラットフォーム「WonderMarket」の設計と実装について報告する．本システムは，同一基盤上で「購入者向けカタログ・購買UI」と「販売者向け店舗・受注管理ダッシュボード」を完全分離かつ動的に同期させる二重モード構造を実現している．さらに，ローカルサーバー環境（ポート80）でのJSONファイル永続化と，GitHub Pages等の静的ホスティング環境でのLocalStorageフォールバック機構を自動判別するデュアル実行機構を導入し，高いポータビリティと実用性を両立した．
\end{abstract}

% ==========================================================
% 1. はじめに
% ==========================================================
\section{はじめに}
現代のEコマース（EC）システムは，決済・在庫・注文処理などの多岐にわたる機能を持ち，その開発・運用環境の複雑化が課題となっている．特に初学者教育やプロトタイピング環境においては，過剰な依存関係（重厚なフレームワークや外部DB）を排しつつ，実運用に近い購買・販売サイクルを安全に検証できる軽量なシステムが求められる．

本研究では，Python標準ライブラリのみで動作するRESTfulサーバーと，依存ライブラリゼロのモダンフロントエンドを組み合わせ，\textbf{即時ブラウザ起動}，\textbf{実金銭取引のない安全な疑似決済}，および\textbf{GitHub Pages対応の静的公開}を同時に満たす統合Webアプリケーションを開発した．

% ==========================================================
% 2. システム設計と技術選定
% ==========================================================
\section{システムアーキテクチャ}
本システムのアーキテクチャの根幹は「トリプルフォールバック型データ層」と「状態共有フロントエンド」である．

\subsection{バックエンド: Python標準HTTP/API}
外部ライブラリ（FlaskやFastAPI等）に依存せず，Python 3標準の\texttt{http.server}モジュールを拡張して実装した．
\begin{itemize}
  \item \textbf{ポート80直接待受}: \url{http://localhost} で追加ポート番号不要のシームレスなアクセスを実現．
  \item \textbf{RESTfulエンドポイント}: 商品CRUD（\texttt{/api/products}），注文・在庫減算処理（\texttt{/api/orders}），ステータス変更（PATCH），およびデモ初期化（\texttt{/api/reset}）を提供．
  \item \textbf{JSONファイル永続化}: データベースサーバー不要で，\texttt{data/*.json}に原子性をもって保存．
\end{itemize}

\subsection{フロントエンド: コンポーネント指向Vanilla JS}
ReactやVueを用いず，Web標準のES ModulesとDOM APIでSPAライクな画面遷移・リアクティブUIを構築した．画像切損を防ぐため，全商品アイコンはインラインSVGベクターとして動的描画される．

% ==========================================================
% 3. 主要機能の実装
% ==========================================================
\section{主要機能の実装}

\subsection{購入者向けストアフロント (Buyer Mode)}
一般消費者が直感的に閲覧・購買できるUIを提供する．
\begin{enumerate}
  \item \textbf{動的カタログ}: リアルタイムキーワード検索，カテゴリ抽出，価格・評価順ソートを即時実行．
  \item \textbf{商品詳細モーダル}: 高精細SVGプレビュー，在庫アラート，スペック表示．
  \item \textbf{スライドイン・カート}: 非同期数量変更，10\%消費税・合計額の自動計算．
  \item \textbf{デモチェックアウト}: 配送先入力後，疑似決済をシミュレート．クレジットカード等実取引は発生せず，注文番号を発行して完了画面へ遷移．
\end{enumerate}

\subsection{販売者向けダッシュボード (Seller Mode)}
店舗運営者が受注と在庫を一括管理するコンソールである．
\begin{enumerate}
  \item \textbf{KPIサマリー}: 総売上高，注文件数，出品数，在庫僅少警告（5点以下）をカード表示．
  \item \textbf{商品出品・在庫管理}: 新規出品モーダル，商品編集/削除，テーブル上での直接在庫増減（$+$ / $-$）．
  \item \textbf{受注・発送追跡}: 顧客情報と注文明細の確認，「受付$\rightarrow$準備中$\rightarrow$発送済$\rightarrow$完了」のステータス更新．
  \item \textbf{売上アナリティクス}: 週間売上推移バーチャート，カテゴリ別売上比率プログレスバーを自動集計描画．
\end{enumerate}

% ==========================================================
% 4. ハイブリッド同期とフォールバック
% ==========================================================
\section{ハイブリッド同期機構}
本システムのデータ層（\texttt{store.js}）は，起動時に自動環境判別を行う．
\begin{itemize}
  \item \textbf{ローカル実行時}: \texttt{/api/products} への通信に成功した場合，Pythonサーバーと同期し，他端末や再起動後もデータが保持される．
  \item \textbf{静的ホスティング時}: GitHub Pages等の静的サーバーや直接ファイルオープン時は，自動的にブラウザ内\texttt{localStorage}へフォールバックする．
\end{itemize}
この設計により，同一コードベースのままGitHubリポジトリとして世界中にWeb公開可能となった．

% ==========================================================
% 5. むすび
% ==========================================================
\section{おわりに}
本稿では，軽量・高ポータブルなECアプリ「WonderMarket」の実装について報告した．ポート80での即時起動，購入者・販売者間の完全なデータ連動，GitHub Pages公開適合性を実証した．今後の課題として，ユーザー認証機構の追加やPWA化が挙げられる．

\end{document}
